\documentclass[12pt]{article}

\usepackage[utf8]{inputenc}
\usepackage[T1]{fontenc}
\usepackage[spanish]{babel}
\usepackage{amsmath}
\usepackage{amssymb}
\usepackage{tikz}
\usepackage{xcolor}
\usepackage{geometry}
\usepackage{enumitem}

% ------------------------------------------------
% MÁRGENES
% ------------------------------------------------
\geometry{
    a4paper,
    left=1.5cm,
    right=1.5cm,
    top=1.3cm,
    bottom=1.5cm
}

% Sin sangría al comenzar párrafos
\setlength{\parindent}{0pt}

% Elimina encabezado y número de página
\pagestyle{empty}

\begin{document}

\noindent
\textbf{3. \textcolor{red}{(20 pts)}} 
En el plano $xy$ se tiene un arco de circunferencia de radio $R$, centrado en el origen $O$ y posee una densidad lineal de carga no uniforme dada por $\lambda(\theta)=\lambda_1\cos\theta$, donde $\lambda_1$ es una constante conocida. A partir del punto del arco ubicado sobre el eje $x$ negativo se extiende una barra recta de longitud $2R$, con densidad lineal 
de carga uniforme $\lambda_0$. Ambas distribuciones de carga se muestran en la figura.

\vspace{0.25cm}

\noindent
\begin{minipage}[t]{0.58\textwidth}
\vspace{0pt}

\textbf{Determine:}

\vspace{0.25cm}

\begin{enumerate}[
    label=\alph*),
    leftmargin=0.55cm,
    labelsep=0.15cm,
    itemsep=0.55cm,
    topsep=0pt
]

    \item El potencial eléctrico total en el origen $O$ debido a 
    ambas distribuciones de carga. 
    \textcolor{red}{\textbf{(8 pts)}}

    \item El valor de $\lambda_0$, en función de $\lambda_1$, para 
    que el potencial eléctrico total en el origen sea nulo. 
    \textcolor{red}{\textbf{(6 pts)}}

    \item Utilizando el valor de $\lambda_0$ obtenido en el apartado 
    anterior, determine la carga eléctrica total producida por ambas distribuciones. 
    \textcolor{red}{\textbf{(6 pts)}}

\end{enumerate}

\end{minipage}
\hfill
%
\begin{minipage}[t]{0.40\textwidth}
\vspace{-0.45cm}
\centering

\begin{tikzpicture}[scale=1.12]

    % Radio gráfico
    \def\R{1.45}

    % ------------------------------------------------
    % EJES
    % ------------------------------------------------
    \draw[->,thick]
        (-5.15,0) -- (1.45,0)
        node[right] {$x$};

    \draw[->,thick]
        (0,-2.05) -- (0,2.15)
        node[above] {$y$};

    % Origen
    \fill (0,0) circle (1.2pt);
    \node[below right] at (0,0) {$O$};

    % ------------------------------------------------
    % ARCO
    % theta = 3pi/4 hasta 5pi/4
    % ------------------------------------------------
    \draw[line width=1.6pt]
        (135:\R)
        arc[
            start angle=135,
            end angle=225,
            radius=\R
        ];

    % Radios a los extremos del arco
    \draw[densely dashed]
        (0,0) -- (135:\R);

    \draw[densely dashed]
        (0,0) -- (225:\R);

    % ------------------------------------------------
    % ÁNGULOS DE 45 GRADOS
    % ------------------------------------------------



    \node at (-0.58,0.25)
        {$45^\circ$};

  
    \node at (-0.58,-0.25)
        {$45^\circ$};

    % ------------------------------------------------
    % DENSIDAD DEL ARCO
    % ------------------------------------------------
    \node[left]
        at (0.9,1.25)
        {$\lambda(\theta)=\lambda_1\cos\theta$};

    % Radio R
    \node[right]
        at (-0.78,-0.95)
        {$R$};

    % ------------------------------------------------
    % BARRA
    % Desde x=-3R hasta x=-R
    % ------------------------------------------------
    \draw[line width=2.5pt]
        ({-3*\R},0) -- (-\R,0);

    % Extremos de la barra
    \draw
        (-\R,0.11) -- (-\R,-0.11);

    \draw
        ({-3*\R},0.11) -- ({-3*\R},-0.11);

    % Densidad de la barra
    \node[above]
        at ({-2*\R},0.20)
        {$\lambda_0$};

    % ------------------------------------------------
    % LONGITUD 2R
    % ------------------------------------------------
    \draw[<->]
        ({-3*\R},-0.48) -- (-\R,-0.48);

    \node[
        fill=white,
        inner sep=1pt
    ]
        at ({-2*\R},-0.48)
        {$2R$};

\end{tikzpicture}

\end{minipage}

\end{document}